\subsection{A Question-Driven Model Story}

\begin{frame}{Year prediction: diagnose the bottleneck, then upgrade}

\small

\begin{center}
\begin{tikzpicture}[
	storybox/.style={
		rounded corners=2pt,
		draw=structbgcolor,
		fill=structbgcolor!10,
		align=center,
		minimum height=1.05cm,
		text width=2.05cm,
		font=\scriptsize
	},
	arrow/.style={->,>=stealth,thick,draw=structbgcolor!75}
]
	\node[storybox] (ridge) at (0,1.3)
		{\textbf{T90 Ridge}\\transparent baseline};
	\node[storybox] (rff) at (3.0,1.3)
		{\textbf{RFF + Ridge}\\test smooth nonlinearity};
	\node[storybox] (tree) at (6.0,1.3)
		{\textbf{LightGBM T90}\\change function family};
	\node[storybox] (audio) at (9.0,1.3)
		{\textbf{594 audio}\\restore information};
	\node[storybox] (ordinal) at (7.5,-1.0)
		{\textbf{Ordinal-MoE}\\target MAE and order};
	\node[storybox] (fused) at (4.5,-1.0)
		{\textbf{Audio + metadata}\\add context};
	\node[storybox] (ensemble) at (1.5,-1.0)
		{\textbf{Ensemble}\\combine residual errors};

	\draw[arrow] (ridge) -- (rff);
	\draw[arrow] (rff) -- (tree);
	\draw[arrow] (tree) -- (audio);
	\draw[arrow,rounded corners=5pt]
		(audio.south) -- (9.0,0.05) -- (ordinal.east);
	\draw[arrow] (ordinal) -- (fused);
	\draw[arrow] (fused) -- (ensemble);
\end{tikzpicture}
\end{center}

\begin{columns}[c,onlytextwidth]
	\column{.32\textwidth}\centering
	{\scriptsize\bfseries\color{black!55}EXECUTION}\par
	{\large\bfseries\color{structbgcolor}Spark}
	\column{.32\textwidth}\centering
	{\scriptsize\bfseries\color{black!55}FUNCTION FAMILY}\par
	{\large\bfseries\color{structbgcolor}Linear $\rightarrow$ Trees}
	\column{.32\textwidth}\centering
	{\scriptsize\bfseries\color{black!55}EVIDENCE}\par
	{\large\bfseries\color{structbgcolor}Audio + Context}
\end{columns}

\note{
Introduce the sequence as a chain of hypotheses, not a leaderboard.
Each arrow changes one main assumption and is justified by the preceding result.
}

\end{frame}

\begin{frame}{Problem and evidence contract}

\centering
\begin{tikzpicture}[
	split/.style={rounded corners=2pt,draw=structbgcolor,line width=1.1pt,
		minimum width=3.15cm,minimum height=1.45cm,align=center},
	flow/.style={->,>=stealth,line width=1pt,draw=black!35}
]
	\node[split,fill=structbgcolor!12] (train) at (-4.0,1.25)
		{\scriptsize\bfseries TRAIN\\[3pt]\Large\bfseries 420,013};
	\node[split] (val) at (0,1.25)
		{\scriptsize\bfseries VALIDATION\\[3pt]\Large\bfseries 46,127};
	\node[split] (test) at (4.0,1.25)
		{\scriptsize\bfseries TEST\\[3pt]\Large\bfseries 49,436};
	\draw[flow] (train) -- (val);
	\draw[flow] (val) -- (test);
	\node[font=\large\bfseries,text=structbgcolor] at (0,-.35)
		{Artist-disjoint: one artist, one split};
\end{tikzpicture}

\vspace{.35cm}

\begin{columns}[c,onlytextwidth]
	\column{.48\textwidth}\centering
	{\scriptsize\bfseries\color{black!55}MAE}\par
	{\Large\bfseries Typical error}
	\column{.48\textwidth}\centering
	{\scriptsize\bfseries\color{black!55}RMSE}\par
	{\Large\bfseries Large-miss penalty}
\end{columns}

\note{
Explain that artist isolation controls identity leakage, but later target-derived
metadata still needs its own train-only and self-exclusion rules.
}

\end{frame}

\begin{frame}{Sanity baselines: metrics prefer different constants}

\centering
\begin{tikzpicture}[
	card/.style={rounded corners=2pt,minimum width=5.1cm,minimum height=3.1cm,
		align=center,line width=1.1pt}
]
	\node[card,draw=audioblue,fill=audiofill] at (-2.9,.55) {
		{\scriptsize\bfseries MEAN PREDICTOR}\\[7pt]
		{\LARGE\bfseries 8.1435}\, {\scriptsize MAE}\\[5pt]
		{\Large\bfseries 10.8019}\, {\scriptsize RMSE}
	};
	\node[card,draw=summaryred,fill=summaryfill] at (2.9,.55) {
		{\scriptsize\bfseries MEDIAN PREDICTOR}\\[7pt]
		{\LARGE\bfseries 7.5812}\, {\scriptsize MAE}\\[5pt]
		{\Large\bfseries 11.3558}\, {\scriptsize RMSE}
	};
	\node[font=\large\bfseries,text=structbgcolor] at (0,-2.0)
		{MSD reference reproduced: 8.13 MAE / 10.80 RMSE};
\end{tikzpicture}

\note{
Use the two constants to explain why MAE and RMSE can rank models differently.
This prepares the audience for the later Ordinal-MoE trade-off.
}

\end{frame}

\begin{frame}{T90 Ridge: a distributed reference}

\centering
\begin{tikzpicture}[
	part/.style={rounded corners=2pt,draw=audioblue,fill=audiofill,
		minimum width=2.4cm,minimum height=1.15cm,align=center},
	reduce/.style={rounded corners=2pt,draw=structbgcolor,fill=structbgcolor!14,
		minimum width=2.5cm,minimum height=1.25cm,align=center},
	flow/.style={->,>=stealth,line width=1.1pt,draw=black!45}
]
	\node[part] (p1) at (-4.3,1.2) {Partition 1\\local gradient};
	\node[part] (p2) at (-4.3,-.3) {Partition 2\\local gradient};
	\node[part] (p3) at (-4.3,-1.8) {Partition $P$\\local gradient};
	\node[reduce] (sum) at (0,-.3) {treeReduce\\gradient sum};
	\node[reduce] (update) at (4.0,-.3) {global Ridge\\update};
	\draw[flow] (p1) -- (sum);
	\draw[flow] (p2) -- (sum);
	\draw[flow] (p3) -- (sum);
	\draw[flow] (sum) -- (update);
\end{tikzpicture}

\vspace{.3cm}

\begin{columns}[c,onlytextwidth]
	\column{.45\textwidth}\centering
	{\scriptsize\bfseries\color{black!55}PAPER-ALIGNED INPUT}\par\vspace{.10cm}
	{\LARGE\bfseries\color{structbgcolor}90}\par
	{\scriptsize timbre means + covariances}
	\column{.45\textwidth}\centering
	{\scriptsize\bfseries\color{black!55}TEST MAE / RMSE}\par
	{\LARGE\bfseries\color{structbgcolor}6.7465 / 9.4074}
\end{columns}

\note{
Clarify that this is paper-aligned, not an exact reproduction of the original
Vowpal Wabbit configuration. L2 stabilizes correlated timbre statistics.
}

\end{frame}

\begin{frame}{Mini-batch Ridge: same model, less gradient work}

\centering
\includegraphics[width=.94\textwidth]{ridge_batching_convergence.pdf}

\vspace{.15cm}

\begin{columns}[c,onlytextwidth]
	\column{.32\textwidth}
	\centering
	{\scriptsize\bfseries FULL BATCH}\par
	{\Large\bfseries 754.99 s}\par
	{\scriptsize 100 equivalent passes}
	\column{.32\textwidth}
	\centering
	{\scriptsize\bfseries 25\% MINI-BATCH}\par
	{\Large\bfseries 577.82 s}\par
	{\scriptsize 25 equivalent passes}
	\column{.32\textwidth}
	\centering
	{\scriptsize\bfseries 10\% MINI-BATCH}\par
	{\Large\bfseries 579.12 s}\par
	{\scriptsize 10 equivalent passes}
\end{columns}

\vspace{.12cm}{\small Test quality remains effectively unchanged: $\mathrm{RMSE}\approx9.407$.}

\note{
Point to the right panel: both mini-batches reach the full-batch final
validation MAE near update 95, after about 23.7 and 9.5 equivalent passes.
Final metrics are recomputed on each complete split.
}

\end{frame}

\begin{frame}{Mini-batching exposes Spark's fixed costs}

\begin{columns}[c,onlytextwidth]
	\column{.76\textwidth}
	\includegraphics[width=\linewidth]{ridge_batching_tradeoff.pdf}
	\column{.21\textwidth}
	\centering
	{\scriptsize\bfseries TOTAL TIME}\par\vspace{.10cm}
	{\LARGE\bfseries\color{structbgcolor}76.5\%}\par
	{\scriptsize with 25\% data passes}
	\vspace{.55cm}
	{\scriptsize\bfseries EXTRA GAIN AT 10\%}\par
	{\LARGE\bfseries\color{summaryred}$\approx 0$}\par
	{\scriptsize fixed costs dominate}
\end{columns}

\medskip

\centering
{\large\bfseries Less gradient work $\neq$ proportional wall-clock speedup}\par
{\scriptsize Data traversal, job startup, scheduling, and validation set the floor.}
\note{
Do not claim cluster-wide linear speedup from one local[4] timing experiment.
The evidence isolates useful gradient work from framework and pipeline overhead.
}

\end{frame}
